\documentclass[12pt]{article}

% ============================================================
% PACKAGES
% ============================================================

\usepackage{amsmath}      % Mathematical environments: align, equation, etc.
\usepackage{amsthm}       % Theorem, definition, proof environments
\usepackage{amsfonts}     % Mathematical fonts such as \mathbb
\usepackage{amssymb}      % Additional mathematical symbols
\usepackage{mathtools}    % Extensions to amsmath
\usepackage{geometry}     % Page margins
\usepackage{enumitem}     % Customize enumerate/itemize
\usepackage{mdframed}     % Framed boxes
\usepackage{tikz}         % Draw diagrams
\usepackage{hyperref}     % Clickable references and links

\geometry{
    letterpaper,
    margin=0.8in
}

\linespread{1.2}


% ============================================================
% NUMBER SYSTEMS
% ============================================================
% Instead of writing \mathbb{R} every time, use \RR.

\newcommand{\RR}{\mathbb{R}}   % Real numbers
\newcommand{\CC}{\mathbb{C}}   % Complex numbers
\newcommand{\NN}{\mathbb{N}}   % Natural numbers
\newcommand{\ZZ}{\mathbb{Z}}   % Integers
\newcommand{\QQ}{\mathbb{Q}}   % Rational numbers


% ============================================================
% COMMON MATHEMATICAL NOTATION
% ============================================================
% These commands are useful for frequently used notation.

\newcommand{\set}[1]{\left\{#1\right\}}
% Example:
%   \set{x \in \RR : x > 0}
% produces:
%   { x in R : x > 0 }

\newcommand{\abs}[1]{\left|#1\right|}
% Example:
%   \abs{x-y}
% produces:
%   |x-y|

\newcommand{\norm}[1]{\left\|#1\right\|}
% Example:
%   \norm{x}
% produces:
%   ||x||

\newcommand{\inner}[2]{\left\langle #1,#2\right\rangle}
% Example:
%   \inner{x}{y}
% produces:
%   <x,y>

\newcommand{\parens}[1]{\left(#1\right)}
% Example:
%   \parens{x+y}

\newcommand{\brac}[1]{\left[#1\right]}
% Example:
%   \brac{a,b}

\newcommand{\angles}[1]{\left\langle #1\right\rangle}
% Example:
%   \angles{x,y}


% ============================================================
% THEOREM / THEORY COMPONENTS
% ============================================================

% "definition" style is upright text.
% Use this for definitions, examples, and remarks.

\theoremstyle{definition}

% ------------------------------------------------------------
% DEFINITION
% ------------------------------------------------------------
% Use when introducing a formal mathematical definition.
%
% Usage:
%
% \begin{definition}[Metric]
% A metric on a set X is a function ...
% \end{definition}
%
% [Metric] is optional and gives the definition a name.

\newtheorem{definition}{Definition}[section]


% ------------------------------------------------------------
% EXAMPLE
% ------------------------------------------------------------
% Use to give a concrete example of a definition/theorem.
%
% Usage:
%
% \begin{example}
% Consider X = \RR with the usual distance.
% \end{example}

\newtheorem{example}[definition]{Example}


% ------------------------------------------------------------
% REMARK
% ------------------------------------------------------------
% Use for observations, comments, or small additional facts.
% It is NOT intended for major mathematical results.
%
% Usage:
%
% \begin{remark}
% This definition does not require X to be a subset of \RR^n.
% \end{remark}

\newtheorem{remark}[definition]{Remark}


% ============================================================
% THEOREM-LIKE COMPONENTS
% ============================================================

% "plain" style makes the mathematical statements italic.

\theoremstyle{plain}


% ------------------------------------------------------------
% THEOREM
% ------------------------------------------------------------
% Use for an important/general mathematical result.
%
% Usage:
%
% \begin{theorem}[Cauchy--Schwarz Inequality]
% For all x,y \in \RR^n,
% \[
% |\inner{x}{y}| \leq \norm{x}\norm{y}.
% \]
% \end{theorem}

\newtheorem{theorem}[definition]{Theorem}


% ------------------------------------------------------------
% PROPOSITION
% ------------------------------------------------------------
% Use for a useful mathematical result that is usually less
% central than a theorem.
%
% Usage:
%
% \begin{proposition}
% Every finite-dimensional subspace of \RR^n is closed.
% \end{proposition}

\newtheorem{proposition}[definition]{Proposition}


% ------------------------------------------------------------
% LEMMA
% ------------------------------------------------------------
% Use for an auxiliary result whose main purpose is to help
% prove another theorem/proposition.
%
% Usage:
%
% \begin{lemma}
% ...
% \end{lemma}

\newtheorem{lemma}[definition]{Lemma}




% ------------------------------------------------------------
% COROLLARY
% ------------------------------------------------------------
% Use for a result that follows relatively directly from a
% theorem/proposition that was just established.
%
% Usage:
%
% \begin{corollary}
% ...
% \end{corollary}

\newtheorem{corollary}[definition]{Corollary}


% ============================================================
% PROOF
% ============================================================
% The proof environment comes from amsthm.
%
% Usage:
%
% \begin{proof}
% Suppose that ...
%
% Therefore, ...
% \end{proof}
%
% LaTeX automatically adds the QED symbol at the end.


% ============================================================
% CUSTOM BOX: IDEA
% ============================================================
% Use this when you want to record intuition or the main idea
% behind a mathematical concept.
%
% Usage:
%
% \begin{idea}
% \textbf{Key idea.}
% A metric gives us a notion of distance...
% \end{idea}

\newmdenv[
    linewidth=1pt,
    skipabove=10pt,
    skipbelow=10pt
]{idea}


% ============================================================
% CUSTOM BOX: WARNING
% ============================================================
% Use this for common mistakes, distinctions, or things that
% are easy to confuse.
%
% Usage:
%
% \begin{warning}
% Do not confuse continuity with uniform continuity.
% \end{warning}

\newmdenv[
    linewidth=1pt,
    skipabove=10pt,
    skipbelow=10pt
]{warning}

% ============================================================
% CUSTOM BOX: NOTE
% ============================================================
% Use this for a short personal note, useful observation,
% connection between concepts, or explanation.
%
% Unlike "Remark", this is intended for informal notes.
%
% Usage:
%
% \begin{note}
% A sequence can be viewed as a function from $\NN$ to $\RR$.
% \end{note}

\newmdenv[
    linewidth=1pt,
    skipabove=10pt,
    skipbelow=10pt
]{note}


% ============================================================
% DOCUMENT
% ============================================================

\begin{document}


% ============================================================
% TITLE
% ============================================================

\begin{center}
    {\Large \textbf{Properties of $\mathbb{Q}$}}\\[0.5em]
    {\large Lecture 2}
\end{center}

% Automatically generate the table of contents.
\tableofcontents

\newpage
\section{Arithmetic on Rational numbers}
\subsection{Well-defined operators}
\begin{definition}[Well-defined-ness]
Operators are called well-defined if their results do not depend on the choice of representative.
\end{definition}
\begin{example}
$\frac{a}{b} + \frac{c}{d} = \frac{a+c}{b+d}$ is not well-defined. We have $\frac{1}{2} + \frac{1}{3} = \frac{2}{5}$ but $\frac{2}{5}$ is not equivalent with $\frac{3}{7} = \frac{2}{4} + \frac{1}{3}$. Thus, we have "something" in $\frac{1}{2}$ plus with $\frac{1}{3}$ has different outcomes.
\end{example}

\begin{example}
$\frac{a}{b} + \frac{c}{d} = \frac{0}{1}$ is well-defined, its result will never depend on the choice of representative. However, this arithmetic makes the arithmetic on $\ZZ$ doesn't work same as before.
\end{example}

\subsection{Addition and Multiplication on Rational numbers}

\begin{definition}[Addition in $\QQ$]
$\frac{a}{b}+\frac{c}{d} = \frac{ad+bc}{bd}$
\end{definition}
\begin{note}
To check the well-defined-ness of the addition above, we have to check if:
$(a,b) \sim (a', b')$ and $(c,d) \sim (c', d')$ then $(ad+bc,bd) \sim (a'd'+b'c', bd)$

\textit{Proof}:

$(ad+bc)b'd' = adb'd'+bcb'd'$

$= a'bdd'+c'dbb'$ (Given $ab'=a'b$ and $cd'=c'd$ because $(a,b)\sim(a',b')$ and $(c,d)\sim(c',d')$)

$= (a'd'+b'c')bd$

$\Rightarrow (ad+bc,bd) \sim (a'd'+b'c', bd)$
\end{note}

\begin{definition}[Multiplication on $\QQ$]
$\frac{a}{b} \frac{c}{d} = \frac{ac}{bd}$
\end{definition}

\begin{note}
Same concept for proving well-defined of multiplication in $\QQ$.
\end{note}

\subsection{Order in Rational numbers}
\begin{definition}[Order]
An order on set $S$ is a relation, denoted as "$<$" satisfying:
\begin{itemize}
    \item \textbf{Trichotomy} Consider any 2 elements $x$ and $y$ in $S$, $x$ and $y$ will fall only to 1 relation in 3 given relations: $x <y$, $x=y$, $y<x$.
    \item \textbf{Transitive} $x<y$ and $y<z$ then $x<z$.
\end{itemize}
We called $S$ is an ordered set, denoted as $(S, <)$.
\end{definition}

\begin{example}
In $\ZZ$, we have $m<n$ if $n-m \in \{1,2,3,4,5,...\}$ (positive).

In $\QQ$, we have $\frac{m}{n}$ is positive if $mn > 0$ (both $m,n$ are positive or both non-positive).

Hence, in $\QQ$, $\frac{a}{b}<\frac{c}{d}$ if $\frac{-a}{b}+\frac{c}{d}$ is positive.
\end{example}

\section{Field}
\begin{definition}
A field is a set that is equipped with 2 operators: addition and multiplication, and satisfied some axioms.
\end{definition}

\begin{note}
There are 11 axioms in total, 5 for addition, 5 for multiplication, 1 for distribution between multiplication and addition. I won't mention in this note.
\end{note}

\begin{note}
It is good to know that $\QQ$ is a field, but $\ZZ$ is not. $\QQ$ is even a ordered field: field with order and the order preserved by field operators ($\times$, $+$):
\begin{itemize}
    \item $y<z \Rightarrow x+y<x+z$.
    \item $y<z$ and $x>0$ then $xy>xz$.
\end{itemize}
\end{note}

\begin{note}
Note that, not every field equipped with order is an ordered field, it must satisfy 2 conditions above.
\end{note}

\section{Why Rational?}
\begin{note}
$\QQ$ helps us solve some equation that can't be solved in $\ZZ$ (i.e. $5x=3$).

But $\QQ$ is not enough. Look at the under theorem.
\end{note}

\begin{theorem}
$x^2 = 2$ has no solution in $\QQ$.
\end{theorem}

\begin{note}
That's why we will construct $\RR$!
\end{note}


\end{document}